\subsection{THE MEAN VALUE THEOREM FOR REAL-VALUED FUNCTIONS}

The following important result is a consequence of the corollary to prop. 1: if one has \(m \leqslant f'(x) \leqslant M\) on \(]a, b[\) , then also \(m \leqslant \frac{f(b) - f(a)}{b - a} \leqslant M\) . In other words, a bound for the derivative of \(f'\) on the whole interval with endpoints \(a, b\) implies the same bound for \(\frac{f(b) - f(a)}{b - a}\) (the ratio of the ``increment'' of the function to the ``increment'' of the variable on the interval). We shall make this fundamental result more precise, and generalize it, in the sequel.

PROPOSITION 2. Let \(f\) be a real function which is finite and continuous on the closed bounded interval \(I = [a, b]\) (where \(a < b\) ) and has a right derivative (finite or not) at all the points of the relative complement in \([a, b)\) of a countable subset \(A\) of this interval. If \(f_d'(x) \geqslant 0\) at every point of \([a, b[\) not belonging to \(A\) , then one has \(f(b) \geqslant f(a)\) ; if, further, \(f_d'(x) > 0\) for at least one point of \([a, b[\) , then \(f(b) > f(a)\) .

Let \(\varepsilon > 0\) be arbitrary, and denote by \((a_n)_{n \geqslant 1}\) a sequence obtained by listing the countable set A. Let J be the set of points \(y \in I\) such that one has

\[
f (x) - f (a) \geqslant - \varepsilon (x - a) - \varepsilon \sum_ {a _ {n} <   x} \frac {1}{2 ^ {n}}\tag{1}
\]

for all \(x\) with \(a \leqslant x \leqslant y\) , the sum in the second term of the right-hand side being taken over all indices \(n\) for which \(a_n < x\) . We shall show that if \(f_d'(x) \geqslant 0\) at every point of \([a, b[\) distinct from the \(a_n\) , then \(J = I\) .

It is clear that J is not empty, since \(a \in J\) ; moreover the definition of this set shows that if \(y \in J\) one has \(x \in J\) for \(a \leqslant x \leqslant y\) , so J is an interval with left-hand endpoint a (Gen.~Top., IV, p.~336, prop. 1); let c be its right-hand endpoint. One has \(c \in J\) ; this is clear if c = a; if not, for every x \textless{} c we have the inequality (1), and a fortiori

\[
f (x) - f (a) \geqslant - \varepsilon (c - a) - \varepsilon \sum_ {a _ {n} <   c} \frac {1}{2 ^ {n}}
\]

from which it follows, on letting \(x\) tend to \(c\) in this inequality (since \(f\) is continuous), that \(c\) satisfies (1).

This being so, we shall see that we must have c = b. Indeed, if one had c \textless{} b, then certainly one would have \(c \notin A\) ; now \(f_{d}^{\prime}(c)\) exists, and since \(f_{d}^{\prime}(c) \geqslant 0\) by hypothesis, there exists a y such that \(c < y \leqslant b\) and such that for \(c \leqslant x \leqslant y\) one has

\[
f (x) - f (c) \geqslant - \varepsilon (x - c)
\]

from which, taking account of (1), where x is replaced by c,

\[
f (x) - f (a) \geqslant - \varepsilon (x - a) - \varepsilon \sum_ {a _ {n} <   c} \frac {1}{2 ^ {n}} \geqslant - \varepsilon (x - a) - \varepsilon \sum_ {a _ {n} <   x} \frac {1}{2 ^ {n}}
\]

which signifies that \(y \in J\) , contradicting the definition of c.~Thus we have \(c = a_{k}\) for some index k; since f is continuous at the point \(a_{k}\) there is a y such that \(c < y \leqslant b\) and such that for \(c < x \leqslant y\) one has

\[
f (x) - f (c) \geqslant - \frac {\varepsilon}{2 ^ {k}}
\]

from which, taking account of (1), where x is replaced by c,

\[
f (x) - f (a) \geqslant - \varepsilon (c - a) - \varepsilon \sum_ {a _ {n} <   x} \frac {1}{2 ^ {n}} \geqslant - \varepsilon (x - a) - \varepsilon \sum_ {a _ {n} <   x} \frac {1}{2 ^ {n}}
\]

which again leads to a contradiction; we thus have c = b, and in consequence

\[
f (b) - f (a) \geqslant - \varepsilon (b - a) - \varepsilon \sum_ {a _ {n} <   b} \frac {1}{2 ^ {n}} \geqslant - \varepsilon (b - a) - \varepsilon .\tag{2}
\]

Since \(\varepsilon > 0\) is arbitrary we deduce from (2) that \(f(b) \geqslant f(a)\) , which demonstrates the first part of the proposition.

We remark now that this result applied to an interval \([x, y]\) where \(a \leqslant x < y \leqslant b\) proves that f is increasing on I; if one had \(f(b) = f(a)\) one could deduce that f is constant on I, and then that \(f_{d}'(x) = 0\) at every point of \([a, b[\) ; the second part follows from this.

COROLLARY. Let \(f\) be a finite continuous real function on \([a, b]\) (where \(a < b\) ) and have a right derivative at all points of the complement in \([a, b[\) of a countable subset A of this interval. For \(f\) to be increasing on I it is necessary and sufficient that \(f_d'(x) \geqslant 0\) at every point of \([a, b[\) that does not belong to A; for \(f\) to be strictly increasing it is necessary and sufficient that that the preceding condition holds, and further that the set of points \(x\) where \(f_d'(x) > 0\) be dense in \([a, b]\) .

Remarks. 1) Prop. 2 remains true when one replaces the interval \([a, b[\) by \(]a, b]\) and the words ``right derivative'' by ``left derivative''.

\begin{enumerate}
\def\labelenumi{\arabic{enumi})}
\setcounter{enumi}{1}
\item
  The hypothesis of continuity on \(f\) on the closed interval I (and not just right continuity \(^{4}\) at every point of \([a, b]\) ) is essential for the validity of prop. 2 (cf.~I, p.~36, exerc. 8).
\item
  The conclusion of prop. 2 is not guaranteed if one merely supposes that the set A of ``exceptional'' points is nowhere dense in I, but not countable (cf.~I, p.~37, exerc. 3).
\end{enumerate}

Prop. 2 entails the following fundamental theorem (which appears to be more general):

THEOREM 1 (mean value theorem). Let \(f\) and \(g\) be two finite continuous real-valued functions defined on a closed bounded interval \(I = [a, b]\) and having a right derivative (finite or not) at all points of the relative complement in {[}a, b{[} of a countable subset of this interval. Suppose further that \(f_d'(x)\) and \(g_r'(x)\) are not simultaneously infinite except at the points of a countable subset of I and that there are finite numbers \(m, M\) such that

\[
m g _ {r} ^ {\prime} (x) \leqslant f _ {d} ^ {\prime} (x) \leqslant \mathrm{Mg} _ {r} ^ {\prime} (x)\tag{3}
\]

except at the points of a countable subset of I (replacing \(\mathbf{M}g_{r}^{\prime}(x)\) (resp. \(mg_{r}^{\prime}(x)\) ) by 0 if M = 0 (resp. m = 0) and \(g_{r}^{\prime}(x) = \pm\infty\) ). Under these conditions one has

\[
m (g (b) - g (a)) <   f (b) - f (a) <   \mathbf {M} (g (b) - f (a))\tag{4}
\]

except when one has \(f(x) = \mathrm{Mg}(x) + k\) , or \(f(x) = mg(x) + k\) (k constant) for all \(x \in \mathbf{I}\) .

It suffices to apply prop. 2 to the functions Mg - f and f - mg, which, under our hypotheses, have a positive right derivative except at the points of a countable subset of I.

Remark. Th. 1 fails if one allows \(f_{d}^{\prime}\) and \(g_{r}^{\prime}\) to be simultaneously infinite on an uncountable subset of I (cf.~I, p.~37, exerc. 3).

COROLLARY. Let \(f\) be a finite continuous function on \([a, b]\) (where \(a < b\) ) and have a right derivative (finite or not) at all points of the relative complement \(\mathbf{B}\) in \([a, b[\) of a countable subset of this interval. If \(m\) and \(M\) are the greatest lower and least upper bounds of \(f_d'\) on \(B\) then one has

\[
m (b - a) <   f (b) - f (a) <   \mathrm{M} (b - a)\tag{5}
\]

if \(f\) is not an affine linear function; if \(f\) is affine linear one has

\[
m = \mathbf {M} = \frac {f (b) - f (a)}{b - a}.
\]

The inequalities (5) are consequences of (4) when \(m\) and \(M\) are finite; the case when one or the other of these numbers is infinite is trivial.

Remark. The inequalities (5) prove that a continuous function cannot have right derivative equal to \(+\infty\) at all points of an interval (cf.~I, p.~38, exerc. 6).

